\section{What Placement Has To Get Right}
\label{sec:evidence}

\Cref{sec:formulation} frames placement as a learnable budget-constrained MDP. This section reports what
we found when we tried to learn it on production KV traces, and it is organized around a defect rather
than a result, because the defect came first and everything else follows from it: the protocol these
studies use --- ours for three revisions --- makes 97\% of the achievable benefit unreachable, so a
policy measured inside it looks flat however well it ranks (\Cref{sec:reachable}). Moving the decision to
each block's arrival recovers the benefit, and shows that the recoverable part is captured by
\emph{deciding early} rather than by \emph{deciding well} (\Cref{sec:arrival}). We then report what a
better ranker buys once the metric is repaired (\Cref{sec:capacity}), where reaction stands
(\Cref{sec:reaction}), and --- at length, because it is the most transferable part --- the protocol,
including the three claims we retracted along the way (\Cref{sec:protocol}).

\subsection{Sources and references}
\label{sec:sources}
We use \textbf{Mooncake}~\cite{qin2024mooncake}, replayed production traces from a general LLM serving
stack (23{,}608 and 12{,}031 requests), and \textbf{kv-cache-tester}, anonymized agentic coding
conversations (200 of 739 fetched). Each request lists the hash ids of the KV blocks it touches, so a
repeated hash is a genuine reuse. Both sources record wall-clock request timestamps. Block ids in
kv-cache-tester are conversation-scoped, so we namespace them before pooling; not doing so invents
cross-conversation reuse.

Three references bound every policy, all billed on the true future: a \textbf{hindsight oracle} that ranks
and places knowing future access counts; a \textbf{frequency heuristic} that assumes the future looks like
the observed prefix; and a \textbf{naive} all-plaintext baseline. ``Headroom'' is
$(\text{cost}_{\text{heur}} - \text{cost}_{\text{oracle}})/\text{cost}_{\text{heur}}$; ``captured'' is the
fraction of that interval a learned policy recovers.

\subsection{The protocol decides what is reachable}
\label{sec:reachable}
Placement studies of this kind fix a \emph{split instant}: observe the trace up to some point, build
features from that prefix, place every block, and bill the result against what the future turns out to
contain. It is the natural protocol, it is what we ran for three revisions of this paper, and it has a
structural defect that no aggregate computed inside it can reveal.

The defect is general, not particular to our setup. Whenever a workload admits new items after the split
--- which is to say, whenever it is not a closed working set --- a split-time policy has no opinion about
those items while the hindsight oracle it is scored against ranks them freely. The reported ``fraction of
headroom captured'' then divides by a quantity that is partly unreachable, and the unreachable share
grows with the arrival rate. Any study reporting that fraction without bounding it is reporting a number
whose denominator it has not characterized.

Feature construction can only describe a block that has been seen. On these traces
\textbf{45.2\% and 46.1\% of blocks have no pre-split access at all}, and a policy therefore has no
opinion about them --- but the hindsight oracle it is scored against ranks them anyway, using knowledge
of a future in which they exist. Those blocks carry \textbf{96.9\% and 97.6\%} of the oracle's total
advantage. The denominator of every ``captured headroom'' number we reported was, in other words,
almost entirely unreachable by construction:

\begin{center}\small
\begin{tabular}{@{}lrr@{}}
\toprule
& conversation & toolagent \\
\midrule
blocks with no pre-split access & 45.2\% & 46.1\% \\
full hindsight gap & 160{,}200 & 162{,}602 \\
\textbf{attainable gap} (clairvoyance on visible blocks) & \textbf{4{,}996} & \textbf{3{,}900} \\
share of the gap a causal policy cannot reach & \textbf{96.9\%} & \textbf{97.6\%} \\
\bottomrule
\end{tabular}
\end{center}

\medskip
\noindent A policy confined to 3\% of the prize will look flat no matter how well it ranks, and that is
what our earlier ``decoupling'' result measured. We retract it (\Cref{sec:protocol}).

\paragraph{The diagnostic is one line.} Compute the oracle's advantage twice --- once with clairvoyance
over every item, once restricted to the items the policy could have had an opinion about --- and take the
ratio. If it is far from one, the protocol is the finding. It costs a single extra placement pass, it
needs no new data, and running it on the first day would have saved this paper three revisions. We ship
it as \texttt{attainable\_ceiling} in the artifact and recommend it as a standard reporting requirement
for placement and admission studies: \emph{state the reachable share alongside any captured fraction}.

\subsection{Deciding on arrival, not at a split}
\label{sec:arrival}
The 97\% lives in blocks that do not exist yet when a split-time policy decides. The fix is not a better
model but an earlier decision: place each block \emph{when it first appears}, from context available at
that instant --- the shared-prefix session's age, request rate and recency, the block's position within
its request, the request's length. No feature reads the future; the model is fitted on the first half of
arrivals in arrival order and places the second half.

\begin{center}\small
\begin{tabular}{@{}llrrr@{}}
\toprule
trace & arm & AUC & of full gap & of attainable \\
\midrule
conversation & arrival-learned  & 0.574 & 82.8\% & 91.6\% \\
             & \textbf{arrival-FCFS} & 0.500 & \textbf{84.2\%} & \textbf{93.2\%} \\
             & arrival-shuffled & 0.500 & 82.6\% & 91.4\% \\
             & ceiling          & 1.000 & 90.4\% & 100.0\% \\
\midrule
toolagent    & arrival-learned  & 0.635 & 47.0\% & 88.9\% \\
             & \textbf{arrival-FCFS} & 0.500 & \textbf{47.3\%} & \textbf{89.5\%} \\
             & arrival-shuffled & 0.500 & 45.8\% & 86.6\% \\
             & ceiling          & 1.000 & 52.8\% & 100.0\% \\
\bottomrule
\end{tabular}
\end{center}

\medskip
\noindent Three readings, in order of how much they change the design question.

\emph{Timing recovers the benefit.} The reachable share rises from 3.1\% to \textbf{90.4\%} on
conversation and 2.4\% to \textbf{52.8\%} on toolagent --- $29\times$ and $22\times$. The headroom was
never absent; the protocol could not reach it.

\emph{Deciding at all captures almost all of it; deciding well adds nothing.} First-come-first-served
admission with a fixed action reaches 93.2\% and 89.5\% of the ceiling. The learned arm, at AUC 0.574
and 0.635, reaches 91.6\% and 88.9\% --- \emph{slightly worse}. We are not claiming the learned arm is
harmed by learning; at these AUCs the ranking is close enough to noise that spreading beliefs across the
per-item cost function's thresholds costs more than the ordering wins.

\emph{Even the ordering barely matters.} Shuffling the arrival order costs 1.8 and 3.0 points. What is
left --- the overwhelming majority --- is attributable to nothing more than making the decision at the
moment the block appears.

The design implication is the one the earlier framing could not reach: on these workloads, build an
\textbf{admission point}, not a predictor. That is also where the caching literature has been pointing.
Baleen~\cite{S12_wong2024baleen} is the one system among those we survey that decides admission rather
than eviction, and it optimizes a downstream system cost rather than miss rate.

\subsection{What a better ranker buys, once the metric is fixed}
\label{sec:capacity}
The capacity ladder of our earlier draft asked whether a stronger model converts into a better placement.
Under the corrected metric --- ties broken by tier speed rather than enum order, and every belief
quantile-matched onto one magnitude multiset so that only the ordering varies --- the answer is that
ranking quality moves the outcome monotonically and never enough to help:

\begin{center}\small
\begin{tabular}{@{}lrrr@{}}
\toprule
rung & AUC (conv / tool) & \% of attainable (conv) & (tool) \\
\midrule
constant        & 0.50 / 0.50 & $-154.6$ & $-825.7$ \\
random          & 0.49 / 0.50 & $-208.2$ & $-838.6$ \\
$\log$(count)   & 0.64 / 0.71 & $-70.9$  & $-75.8$ \\
logistic        & 0.81 / 0.93 & $-18.8$  & $-32.1$ \\
GBDT            & 0.80 / 0.91 & $-23.9$  & $-33.0$ \\
MLP             & 0.82 / 0.93 & $-8.5$   & $-19.3$ \\
\midrule
ceiling         & 1.00        & $100$    & $100$ \\
\bottomrule
\end{tabular}
\end{center}

\medskip
\noindent The damage falls monotonically as AUC rises ($-154 \to -71 \to -19 \to -8.5$), so ranking
quality does convert --- into less harm. Every arm is nonetheless below the frequency baseline, because
forcing an imperfect ranking onto the oracle's magnitude multiset assigns large beliefs to blocks that
will not be reused.

\paragraph{The magnitude convention decides the sign, and we measured by how much.}
We report the table above under a stated convention because it does not survive an unstated one. Running
each rung twice --- once on the scale its own scores naturally produce, once quantile-matched onto a
common multiset, \emph{with the ordering identical in both} --- gives:

\begin{center}\small
\begin{tabular}{@{}llrrr@{}}
\toprule
trace & rung & native scale & common scale & difference \\
\midrule
conversation & logistic & $+75.0$ & $-18.8$ & 93.8 \\
             & $\log$(count) & $-39.2$ & $-70.7$ & 31.5 \\
             & random   & $-81.1$ & $-208.2$ & 127.1 \\
\midrule
toolagent    & logistic & $+84.8$ & $-32.1$ & 117.0 \\
             & $\log$(count) & $-25.2$ & $-59.1$ & 33.9 \\
             & random   & $-630.2$ & $-838.6$ & 208.5 \\
\bottomrule
\end{tabular}
\end{center}
\noindent{\footnotesize Percent of the attainable ceiling. The two columns differ only in the magnitudes
attached to an identical ranking.}

\medskip
\noindent The same policy, with the same ordering, scores $+84.8\%$ or $-32.1\%$ depending on a choice
of scale. The gap runs from \textbf{31.5 to 208.5 points} --- an order of magnitude larger than any
effect this paper or its predecessors reported. An earlier draft asserted that this control ``changes
the captured headroom by about two points''; that number was never computed, and it is wrong by two
orders of magnitude.

We therefore claim neither column. What we do claim is the comparison \emph{within} a column, which is
well defined: on the common scale, ranking quality converts monotonically into less harm. And we claim
the methodological point, which we think is the more useful one: \textbf{``fraction of headroom
captured'' is not a quantity until the magnitude convention is stated}, and no version of this paper,
nor any placement study we are aware of, states it.


\begin{table}[t]
  \centering \small \setlength{\tabcolsep}{4pt}
  \begin{tabular}{@{}lrr@{}}
    \toprule
     & \makecell{Mooncake\\(serving)} & \makecell{kv-cache-tester\\(agentic)} \\
    \midrule
    requests analysed              & 23{,}608 / 12{,}031 & 4{,}215 (40 conv.) \\
    blocks touched exactly once    & 78\%   & 15\% \\
    mean reuse per block           & 1.9    & 24 \\
    blocks per request             & $\sim$17 & $\sim$1{,}900 \\
    block size (tokens)            & 256    & 64 \\
    \midrule
    headroom, heuristic vs.\ oracle & 87\%  & 81--94\% \\
    headroom captured by a learned policy & $+2$\% & not measurable$^\dagger$ \\
    \bottomrule
  \end{tabular}
  \caption{The two sources, at the scales actually analysed. Descriptive statistics are
  sample-dependent: the agentic figures shift with the number of pooled conversations (one-time fraction
  ranges 7.5--21\% over 20--200 conversations), and the captured-headroom range spans a sign reversal
  between 40 and 150 conversations (\Cref{sec:protocol}). We report the spread rather than a point
  because the sampling is not randomized. $^\dagger$Under randomized draws the agentic figure has a 95\%
  spread an order of magnitude wider than any effect (\Cref{sec:sampling}).}
  \label{tab:regimes}
\end{table}


\subsection{A measurement that was never a measurement}
\label{sec:sampling}
We also attempted the obvious follow-up --- fit a richer per-block model on the agentic source and place
with it --- and first obtained what looked like a strong positive: $+66$ to $+88\%$ of the headroom
captured, reproducing across what we took to be two samples. It was not a result, and the way it failed is
worth more than the number would have been.

The two samples were \emph{sorted prefixes} of the conversation corpus, so the smaller was contained in
the larger; and file index correlates strongly with conversation shape (mean requests per conversation
falls from 140 to 37 between files 1--20 and 71--150). Extending the same prefix to 100 and 150
conversations drove the captured headroom to $-9\%$ and $-48\%$ --- a sign reversal. Re-drawing the
samples \emph{at random} rather than as prefixes then showed why: four random draws at 40
conversations land between $-234$ and $+66$ points, and at 20 conversations between $-1207$ and $+78$
(\Cref{tab:kvct}). The spread dwarfs the effect by an order of magnitude. We quote the observed range
rather than a confidence interval on purpose: with four draws a parametric interval rests on a normality
assumption these data cannot support, and \emph{captured} is a ratio whose denominator varies per draw.
An earlier revision of this paper reported $\pm126$ here, computed with the normal quantile against a
standard deviation estimated from four points; the correct $t(3)$ half-width is $\pm205$. That error is
the eleventh entry in \Cref{sec:protocol}'s catalogue and it was ours. Under a correct sampling protocol there was never a
measurement here to report.

We keep this in the paper for two reasons. It is the control that our own protocol
(\Cref{sec:protocol}) lacked, and running it costs minutes; and it bounds what the rest of our numbers
are worth, since every other result in this section is also a single draw of a parameter we chose.

\begin{table}[t]
  \centering \small \setlength{\tabcolsep}{5pt}
  \begin{tabular}{@{}lcc@{}}
    \toprule
    & \makecell{sorted-prefix samples\\(as first reported)} & \makecell{4 random draws\\mean [observed range]} \\
    \midrule
    20 conversations & --- & $-541\%$ $[-1207, +78]$ \\
    40 conversations & $+88\%$ & $-67\%$ $[-234, +66]$ \\
    70 conversations & $+87\%$ & --- \\
    100 conversations & $-9\%$ & --- \\
    150 conversations & $-48\%$ & --- \\
    \bottomrule
  \end{tabular}
  \caption{Headroom captured by a learned placement policy on agentic traces, the same pipeline measured
  two ways. Left: sorted-prefix samples, which are nested and drawn from a corpus whose per-conversation
  statistics vary threefold along the file index --- the numbers we first reported, and their reversal.
  Right: randomized draws at the same sizes, summarized by their observed range --- four draws cannot
  support a parametric interval. The spread is an order of magnitude wider than any effect, so this
  quantity was never measurable under this protocol. We report it as a cautionary measurement.}
  \label{tab:kvct}
\end{table}


\subsection{Where reaction stands}
\label{sec:reaction}
If forecasting does not pay, the alternative is to react. On the full Mooncake traces, reactive tiering
reaches 75--99\% of Belady's hit rate on the toolagent trace and 40--94\% on open-ended conversation
using LRU, LFU or a fill-once cache alone; admitting the adaptive and learned policies below widens the
per-operating-point envelope only to 76--99\% and 42--94\% (\Cref{fig:reactive}).

The interesting question is not the envelope --- that is an oracle over policies, not something anyone
deploys --- but whether any \emph{single} policy attains it, and whether the adaptive designs built for
exactly this situation beat plain rules. We therefore run the family: \textbf{ARC}~\cite{S7_megiddo2003arc},
which shifts a recency/frequency split on ghost hits; \textbf{LeCaR}~\cite{S8_vietri2018lecar} and
\textbf{CACHEUS}~\cite{S9_rodriguez2021cacheus}, which replace ARC's arithmetic update with regret-based
multiplicative weights over two experts; and the learned \textbf{LRB}, which is reported separately below because its
numbers carry an implementation band the others do not.

\begin{center}\small
\begin{tabular}{@{}lrrrrr@{}}
\toprule
trace & tier & best pure & ARC & LeCaR & CACHEUS \\
\midrule
conversation & 2\%  & 39.8 & 29.4 & \textbf{42.3} & 35.9 \\
             & 5\%  & 53.3 & \textbf{57.9} & 56.8 & 53.5 \\
             & 10\% & 76.0 & \textbf{77.5} & 77.2 & 75.3 \\
             & 20\% & \textbf{94.1} & 87.4 & 93.9 & 88.8 \\
\midrule
toolagent    & 2\%  & 74.9 & 70.7 & \textbf{75.8} & 73.9 \\
             & 5\%  & 81.3 & \textbf{83.6} & 83.0 & 82.9 \\
             & 10\% & 91.9 & 92.3 & \textbf{92.9} & 91.8 \\
             & 20\% & 99.3 & 96.3 & \textbf{99.3} & 97.6 \\
\bottomrule
\end{tabular}
\end{center}
\noindent{\footnotesize Percent of Belady's hit rate. ``Best pure'' is the better of LRU, LFU and a
fill-once cache at that operating point.}

\medskip
\noindent Across eight operating points and four designs, \textbf{the best adaptive or learned policy
beats the best pure rule by at most 4.6 points of Belady}, and at two points by nothing at all. Which
design wins also moves: LeCaR at the tightest tier on both traces, ARC in the middle, plain LRU at the
loosest. That is the substantive finding --- not that these policies fail, but that on KV-block workloads
the whole adaptive-replacement family collapses into a band a few points wide, well short of the oracle
at the small tiers where the decision actually matters.

The diagnostics say why, and the reason is specific to KV workloads. Every policy in this family adapts
on a \emph{ghost hit} --- a miss on a block it recently evicted --- whether it updates a split (ARC) or
an expert weight (LeCaR, CACHEUS). But 76--78\% of blocks in these traces are never reused at all, so
most evictions are correct and produce no signal. Measured, ARC receives an adaptation signal on
\textbf{3.5--7.3\% of misses}; the regret-based policies read the same channel, which is why they land
in the same band rather than above it. What signal it does get is not merely sparse but, at
the smallest tier, directionally wrong: at a 2\% tier ARC drives $p$ to 90\% recency on both traces, while
the rule that actually wins there is frequency. A workload whose blocks are mostly singletons starves the
very mechanism adaptive replacement runs on.

\paragraph{And the learned cache policy loses to both --- but only as a range.} The strongest
counter-argument to this paper is LRB~\cite{S6_song2020lrb}, which learns a relaxed-Belady objective from
a far richer feature set than our four trace-local statistics and won on production CDN traces. We
reimplement it: GBDT regression on log time-to-next-request over inter-access deltas and exponentially
decayed counters, sampled eviction, online retraining from a sliding memory window, right-censored
labels.

Our first version of it measured something else entirely. The training buffer was capped by truncating to
its tail; labelled rows arrive one at a time as blocks are re-accessed while censored rows arrive in one
contiguous burst per checkpoint, so from the third retrain onward every training label was identical, the
regressor returned a constant, and eviction degenerated to a uniform draw over the sampled candidates.
A random-replacement control reproduces that column to within 1.5 points at all eight operating points.
We report the repaired numbers and, because two defensible repairs disagree, we report them as a band:

\begin{center}\small
\begin{tabular}{@{}lrrrr@{}}
\toprule
trace & tier & LRB (repaired) & ARC & below ARC? \\
\midrule
conversation & 2\%  & 27.6--29.2 & 29.4 & yes \\
             & 5\%  & 30.9--32.2 & 57.9 & yes \\
             & 10\% & 45.0--45.9 & 77.5 & yes \\
             & 20\% & 75.8--76.5 & 87.4 & yes \\
\midrule
toolagent    & 2\%  & 68.8--69.0 & 70.7 & yes \\
             & 5\%  & 71.1--71.4 & 83.6 & yes \\
             & 10\% & 80.9--83.0 & 92.3 & yes \\
             & 20\% & 94.6--94.9 & 96.2 & yes \\
\bottomrule
\end{tabular}
\end{center}
\noindent{\footnotesize Percent of Belady. The band spans two repairs of the truncation bug: subsampling
all accumulated rows, and dropping rows older than the memory window first. We adopt the second, which is
what a ``sliding memory window'' means.}

\medskip
\noindent \textbf{The band understates the uncertainty.} An independent reimplementation of the same
repair, produced during review of this paper, reports $\approx$63\% at conversation/10\% where ours gives
45.9\%. We cannot reconcile a 17-point difference without that code. So the honest reading is that our
LRB \emph{point} estimates are not citable to better than about twenty points --- wider than several
differences this paper's earlier drafts drew conclusions from. What survives that band is the ordering:
\textbf{LRB is below ARC at all eight operating points under every repair we or the reviewer produced,
including the 63\% reading.} That is the claim we make; the numbers are context for it.

The mechanism appears to be the one that limits ARC, seen from the other side. LRB learns from labels the
trace supplies when a sampled block is requested again; a block that never returns yields a
\emph{right-censored} row, which says only ``not within the window'' and nothing about when. At the
default window \textbf{77--83\% of LRB's training rows are censored}, and 86--89\% at a shorter one. A
population dominated by singletons starves the ghost-hit channel that adaptive replacement runs on and
the regression target that learned replacement runs on, for the same underlying reason.

Three caveats bound this. It is our reimplementation of the published mechanism, not the authors'
artifact. KV blocks are uniform-size, so byte miss ratio collapses to object miss ratio and LRB loses one
of its structural advantages over LRU. And the hyperparameter sweep we ran to check that the result was
not a tuning artifact was run on the \emph{broken} implementation, so it tested the robustness of random
eviction; we have not repeated it.


The eviction signal also inverts by source. On the agent traces LRU reaches 0.88--0.95 of Belady at a 5\%
tier while LFU stays near 0.23--0.27; on Mooncake the ordering reverses below the working-set knee. An
eviction policy tuned on one source is a poor default for the other --- and, as both the ARC and the LRB
result show, reaching for a more sophisticated policy does not by itself resolve the choice.

\begin{figure}[t]
  \centering
  \includegraphics[width=\columnwidth]{fig_reactive.pdf}
  \caption{Reactive tiering as a fraction of Belady's hit rate, full Mooncake traces. Bars are a
  per-operating-point best-of-five envelope, an oracle over policies rather than a policy. Diamonds are
  ARC, which meets the envelope at 5\% and 10\% but falls short at both extremes; crosses are the
  LRB-style learned policy, which is below ARC everywhere.}
  \label{fig:reactive}
\end{figure}

\subsection{Placement is not retrieval}
\label{sec:placement-vs-retrieval}
A tiering scheduler decides what to keep hot; a retriever decides what enters the prompt. Our first
attempt to measure the difference on LoCoMo~\cite{P51_maharana2024locomo} resolved nothing: two runs
ranked the arms oppositely, and across nine comparisons in each, not one survived correction for
multiplicity. A power analysis said why --- separating the query-agnostic arms needed 246--710 scored
questions per arm and we had 30 --- and also said the reversal was not small-sample noise but came from
the two runs drawing \emph{different dialogues}, which is \Cref{sec:sampling}'s failure in a different
corpus.

We therefore re-ran it at the size the analysis asked for: all ten released dialogues, 30 held-out
questions each, \textbf{300 scored questions per arm}, LLM-judged. Because every arm answers the same
questions the comparison is paired, so we test it with McNemar on the discordant pairs rather than
Fisher on independent samples, and report a per-dialogue sign test alongside so a result cannot rest on
one dialogue.

\begin{center}\small
\begin{tabular}{@{}rrrrr@{}}
\toprule
budget & demand & random & centrality & retrieval \\
\midrule
10\% & 12.0 & 6.7 & 1.7 & \textbf{51.7} \\
25\% & 16.7 & 11.3 & 6.3 & \textbf{52.7} \\
50\% & 29.3 & 25.3 & 16.3 & \textbf{55.3} \\
\bottomrule
\end{tabular}
\end{center}
\noindent{\footnotesize QA accuracy (\%), 300 questions per cell. The first three arms are
query-agnostic placement; retrieval sees the question and is therefore better informed by construction.}

\medskip
\noindent \textbf{Six of the nine comparisons now survive Bonferroni correction}, where none did before.
That is three distinct claims, each replicated at three nested budgets over the same 300 questions, not
six independent findings --- and Bonferroni is the right correction precisely because those tests are
dependent, since it controls the family-wise error rate under arbitrary dependence. Holm on the same
nine $p$-values rejects the same six.

\emph{Retrieval dominates placement, decisively.} It leads at every budget, by a factor of four at the
tightest one, with $p < 10^{-15}$ and the same direction in \textbf{10 of 10 dialogues} (sign test
$p = 0.002$). What was previously argued from scoping is now measured: a query-aware retriever and a
query-agnostic placement policy are not competing for the same job, and a placement signal evaluated as
if it were a retrieval signal will always look bad.

\emph{The placement signal is nonetheless real within its own regime.} Ranking turns by evidence demand
beats embedding centrality --- the fair same-regime baseline --- at every budget ($p \le
1.2\times10^{-4}$). At the dialogue level it leads in 8, 10 and 9 dialogues of 10 (sign test $p = 0.008$,
$0.002$, $0.022$); two of those three do not clear the corrected threshold on their own, so the
question-level test is doing the work here. Query-agnostic placement is not noise; it is simply solving a different and
strictly harder problem than retrieval, with no access to the query.

\emph{One comparison is still unresolved, and we say by how much.} Demand versus random selection points
the right way at all three budgets but reaches only $p = 0.029$, $0.068$ and $0.256$, and at the
dialogue level says nothing at all (5--2, 6--2 and 7--3 dialogues; sign test $p = 0.45$, $0.29$, $0.34$).
At the observed discordance rates, 80\% power would need \textbf{758, 1078 and 2589 questions per arm}
at the corrected threshold we judge significance against --- 486, 667 and 1593 if one sizes at an
uncorrected $0.05$, which is the mistake an earlier revision of this paper made. Against the 300 we ran,
the shortfall is between $2.5\times$ and $8.6\times$. We report this as open with its sample size
attached rather than as a null, and note that a required $n$ estimated from the same data whose null it
describes is itself a noisy quantity: across seeds it moves by roughly 10\%.

The reading we take from this is scoping, now with evidence behind it: placement governs latency and
cost, retrieval governs which facts reach the model, and the two compose rather than compete.


\subsection{The representation axis, and what it says about our own formulation}
\label{sec:representation}
\Cref{sec:formulation} decides representation jointly with tier because the two are argued to be
coupled: a compressed item is cheap to serve but may cost accuracy, so the right form depends on the
budget. Every experiment above holds representation at plaintext, which means that argument --- the
reason the joint decision exists --- had never been tested. We test it here, on LoCoMo, with the same
fact ids, questions and gold answers under three representations that differ only in text and token
cost: the turn verbatim, an LLM summary, and an aggressive key--value extract. All three are
\emph{textual}; the vector, graph, latent and parameter forms of \Cref{sec:formulation} are not
instantiated anywhere in this work, so what follows probes the axis rather than exercising it. Budgets are a fraction of
the \emph{plaintext} total for all three, so a compressed form is not quietly handed a larger allowance.

\begin{center}\small
\begin{tabular}{@{}rrrr@{}}
\toprule
budget & plaintext & summary & extract \\
\midrule
5\%  & 5.0  & \textbf{11.0} & 5.0  \\
10\% & 10.0 & \textbf{15.0} & 7.0  \\
25\% & 19.0 & \textbf{26.0} & 21.0 \\
50\% & 32.0 & \textbf{38.0} & 22.0 \\
\bottomrule
\end{tabular}
\end{center}
\noindent{\footnotesize QA accuracy (\%), 100 questions per cell, LLM-judged. Mean compression ratio:
summary $0.58\times$, extract $0.19\times$.}

\medskip
\noindent Two things are visible and one is not. The trade-off the formulation posits is \emph{visible}:
compress far enough and utility appears to collapse, with the $0.19\times$ extract ten points behind
plaintext at the loosest budget --- though that gap, like everything else here, clears no significance
threshold. But the optimum \emph{does not move}. Summary leads at all four budgets, with the
discordant pairs favouring it every time (6--0, 5--0, 11--4, 19--13). If that holds, representation on
this workload need not be re-decided per operating point --- a weaker coupling than we assumed when we
wrote \Cref{sec:formulation}.

What is not visible is significance. \textbf{None of the eight comparisons survives correction} (best
raw $p = 0.031$ against a 0.00625 threshold); at the observed discordance rates, resolving
summary-versus-plaintext needs 188--1138 questions per cell at the corrected threshold (129--718 at an
uncorrected $0.05$) against the 100 we ran. By the rule we impose
on ourselves in \Cref{sec:protocol}, this is a direction with a sample size attached, not a result. We
report it because it is the first time the axis has been exercised at all, and because the direction
points against our own framing rather than for it.


\subsection{What the measurement required}
\label{sec:protocol}
Every positive result we obtained before the final protocol was an artifact. We report the mechanisms
because they generalize, and because the last three were found only by adversarial review rather than by
our own checks.

\paragraph{Reference-frame mismatches (six).} Two quantities that must share a frame --- the same split
instant, ordering, tie rule, or question set --- were silently defined on different ones: a demand signal
computed from the questions being scored; policy features derived from whole-trace access counts; an
oracle that inherited the heuristic's ordering and so showed zero headroom; a train/test split taken over
access-event indices rather than time, which on real traces falls past the end of the trace; a ranking
metric that broke ties by label, inflating a raw count feature from AUC 0.51 to 0.94; and features built
at one split instant while costs were billed at another.

\paragraph{Retracted claims (three).} Each survived our own checks and was found later --- the last two
by adversarial review, after the catalogue below had already been written. \textbf{(i)} A headline
``+66--88\% of headroom captured on agentic traces'' that reversed sign at a larger sample and that
randomized draws showed had never been measurable (\Cref{sec:sampling}). \textbf{(ii)} A
\emph{decoupling} claim --- prediction improves while placement does not --- that was an artifact of
normalizing against an oracle ranking blocks a causal policy cannot see (\Cref{sec:reachable});
normalized against what is attainable, ranking quality converts monotonically. \textbf{(iii)} A reported
LRB column that measured \emph{random eviction}: the training buffer was truncated to its tail, censored
rows arrive in one contiguous burst per checkpoint, and from the third retrain onward every training
label was identical, so the regressor returned a constant and eviction degenerated to a uniform draw.
Two independent repairs of that bug disagree by 19 points, so we withhold the column rather than publish
either.

The common shape is worth naming: in all three the \emph{aggregate looked reasonable}. A flat captured
fraction, a plausible interval, a learned policy landing below ARC --- each is what the hypothesis
predicted, and none of our checks compared the quantity against what it was arithmetically able to be.

\paragraph{Free-parameter extrapolations (three).} Each was a curve read from one point, in a parameter
fixed for convenience. A reactive claim taken at a single cache size reversed under a sweep. A headroom
taken at 6\% of the trace moved six points by full length (81.0\% to 87.2\%), and the placement
benefit read off it peaked mid-sweep rather than at the longest trace. And a learned-placement result taken at 40 and
70 conversations --- \emph{nested prefixes} of a corpus whose per-conversation statistics vary threefold
along the file index --- reversed sign at 150, from $+88\%$ to $-48\%$, and took a headline claim with it.

\paragraph{Believing a report over an observation (one).} An edit script targeted a heading an earlier
revision had removed, matched nothing, and printed success anyway.

\paragraph{The controls.} Assert that the oracle beats the heuristic it is compared against; keep a
no-learning baseline of the same functional form as the learned one; check that no single feature
outperforms the fitted model; before concluding that a signal is unlearnable, vary model capacity while
holding the protocol fixed, so ``the model was too weak'' is excluded rather than assumed
(\Cref{sec:capacity}); compare reference frames explicitly rather than trusting that two code paths
agree; verify an edit landed rather than that it reported success; sweep every parameter chosen for convenience
before stating a result that depends on it --- the one we lacked; and, before reporting any null, compute
the sample size that would have resolved it, because ``we could not tell'' and ``there is nothing there''
are different claims and only the first is ours to make. The harness ships with all of these wired in,
and a checker that re-derives every number in this paper from its artifact and fails if the two have
drifted --- which caught four rounding drifts and one claim that had quietly mixed two hyperparameter
settings.

\subsection{Limitations}
\label{sec:limits}
No real vLLM+LMCache deployment was measured, so any cost-expressed quantity should be read as a range.
We tightened what that range means. Our first sensitivity analysis scaled every tier latency by a common
factor, which moves the absolute scale and leaves the \emph{ratios} between tiers untouched --- and the
ratios are what a placement policy responds to, so it never tested the parameter that mattered. We
therefore derive the table instead: KV bytes per block follow exactly from a model's attention shape
($2 \times \text{layers} \times \text{kv\_heads} \times \text{head\_dim} \times \text{dtype}$ per
token), and per-tier latency is that divided by a declared device-class bandwidth plus a fixed overhead.
Sweeping three attention shapes against five device classes puts headroom at \textbf{64--88\%}, which
contains the 81\% our hand-written table gives --- the conclusion survives, but two things about the
table do not. Its GPU:CPU latency ratio is 1:10 where the derivation gives 1:80 to 1:250, an order of
magnitude in the ratio the policy is most sensitive to; and \textbf{6 of the 15 device classes reorder
the hierarchy} so that latency order and capacity order disagree, which our fixed per-tier capacity table
cannot express at all. Calibrating against a live deployment remains the open gate; this is the honest
intermediate, not a substitute. The
\emph{timing} axis of \Cref{sec:formulation} is not exercised at all, and the \emph{representation} axis
only as the underpowered pilot of \Cref{sec:representation}; neither is touched by the KV traces. Both
sources are KV-block hash streams from transformer prefix caching, so they are independent in provenance
but not in mechanism; with one exemplar per putative regime we cannot separate workload character from
vendor, date, or block size (256 vs.\ 64 tokens). Sampling is not randomized and no confidence interval is
reported for the placement results, though the two remaining free parameters --- trace length and split
instant --- are now swept rather than fixed. Of the learned and adaptive cache literature we evaluate
ARC~\cite{S7_megiddo2003arc}, LeCaR~\cite{S8_vietri2018lecar}, CACHEUS~\cite{S9_rodriguez2021cacheus} and
a reimplementation of LRB~\cite{S6_song2020lrb}; Hawkeye, Mockingjay and Baleen are positioned in
\Cref{sec:related} but not run, the first two because a KV block hash carries no analogue of the program
counter their predictors are indexed by. Our LRB is the
published mechanism rebuilt from the paper, not the authors' artifact --- as are our LeCaR and CACHEUS, whose
CR-LFU expert we approximate by LFU --- and the uniform block size of both sources removes one of LRB's
advantages over LRU, so the underperformance we report should be read as a statement about this workload
rather than about these designs.
